\subsection{Gaussian Naive Bayes}
\label{sec:gaussian_naive_bayes}

\begin{table*}[htbp]
\caption{Hyperparameter Search Space for Gaussian Naive Bayes Classifier.}
\label{tab:IV}
\centering
\begin{tabular}{llc}
\toprule
Component / Hyperparameter & Evaluated Values & Count \\
\midrule
Feature Scaler & None, MinMaxScaler & 2 \\
Dimensionality Reduction (PCA) & None, 0.50, 0.75 & 3 \\
Variance Smoothing & $\text{logspace}(-9, 2, 40)$ ($1.0 \times 10^{-9}$ to $1.0 \times 10^{2}$) & 40 \\
\midrule
\multicolumn{2}{r}{\textbf{Total Grid Combinations}} & \textbf{240 (1,200 fits)} \\
\bottomrule
\end{tabular}
\end{table*}

Berdasarkan Teorema Bayes, pengklasifikasi GNB mengestimasi probabilitas posterior kelas demografis interseksional dengan asumsi bahwa seluruh dimensi fitur kontinu bersifat independen secara bersyarat \cite{ref35}. Fungsi kerapatan probabilitas distribusi normal memodelkan likelihood fitur kontinu terhadap kelas target sesuai formulasi pada \eqref{eq:6}, di mana $x_i$ menyatakan nilai fitur kontinu ke-$i$, $y$ adalah variabel label kelas, $c$ merupakan kategori kelas target, serta $\mu_{c,i}$ dan $\sigma_{c,i}^2$ masing-masing merepresentasikan rata-rata (mean) dan varians fitur ke-$i$ pada kelas $c$ \cite{ref35}:
\begin{equation}
P(x_i \mid y = c) = \frac{1}{\sqrt{2\pi\sigma_{c,i}^2}} \exp\left(-\frac{(x_i - \mu_{c,i})^2}{2\sigma_{c,i}^2}\right)
\label{eq:6}
\end{equation}

Penambahan parameter penghalusan var\_smoothing menjaga stabilitas komputasi terhadap risiko varians mendekati nol pada ruang representasi laten berdimensi tinggi \cite{ref36}. Prosedur penalaan hyperparameter melalui 5-Fold Stratified Cross-Validation mengevaluasi kombinasi penskalaan data, reduksi dimensionalitas PCA, serta 40 interval var\_smoothing yang dieksplorasi secara logaritmik dari $1.0 \times 10^{-9}$ hingga $1.0 \times 10^{2}$, menghasilkan 240 kombinasi grid dengan total 1,200 proses fitting sebagaimana dipaparkan pada Table~\ref{tab:IV}.
